\documentclass[11pt,a4paper]{article}

\usepackage[utf8]{inputenc}
\usepackage[T1]{fontenc}
\usepackage{amsmath,amssymb,amsfonts}
\usepackage{booktabs}
\usepackage{graphicx}
\usepackage{hyperref}
\usepackage[margin=1in]{geometry}
\usepackage{enumitem}
\usepackage{algorithm,algpseudocode}
\usepackage{xcolor}
\usepackage{multirow}
\usepackage{float}

\title{\Large\textbf{Evaluator Preference Collapse: \\Self-Evaluation Bias in Test-Time Agent Evolution}}
\author{
Liu Zewen\\
\textit{Qilu Institute of Technology, School of Software Engineering}\\
\textit{Tai'an, Shandong, China}\\
\texttt{aidless@github.com}
}
\date{May 2026}

\begin{document}
\maketitle

\begin{abstract}
Test-time reinforcement learning (TTRL) promises LLM agents that improve during inference by self-evaluating outputs and adapting strategies online. We expose a fundamental flaw in this paradigm: \textbf{Evaluator Preference Collapse (EPC)}---when a model serves as both executor and evaluator, strategy weights converge to evaluator preferences rather than task-optimal performance. Across a four-condition controlled experiment (80 LLM calls, DeepSeek-chat + GPT-4o-mini), we measure EPC via a novel Preference Collapse Index (PCI). Self-evaluation (PCI=0.461) exhibits $84\%$ more strategy concentration than ground-truth verification (PCI=0.251). \textbf{Critically, all four evaluator conditions produce mutually inconsistent optimal strategies---there is no evaluator-independent optimum.} Cross-model evaluation (GPT-4o-mini) reduces PCI by $17\%$; cross-role prompting provides a cost-free $2\%$ reduction. We argue that the agent self-improvement literature must adopt multi-evaluator baselines, and provide an open-source experimental framework for reproducible EPC measurement.
\end{abstract}

%===========================================================================
\section{Introduction}
%===========================================================================

The rapid maturation of large language models (LLMs) has catalyzed a new generation of AI agents that plan, reason, and execute multi-step tasks autonomously \cite{yao2023react, wu2023autogen, park2023generative}. A particularly compelling direction is \textit{test-time adaptation}---agents that improve their behavior \textit{during inference} without parameter updates, using self-generated feedback to adjust strategies online \cite{shinn2023reflexion, madaan2023selfrefine, sun2020testtime}.

The core idea is seductively simple: an agent generates an output, evaluates it, and updates its future behavior based on the evaluation signal. This closed loop---generate, evaluate, update---enables continuous improvement without human supervision or offline training \cite{yuan2024selfrewarding, chen2024spin, deepseek2024v3}.

\textbf{We expose a fundamental and previously unreported flaw in this paradigm.} When the same model (or same model family) serves as both the executor proposing solutions and the evaluator judging their quality, the evaluation signal is \textit{not} an objective measure of task performance. Instead, it reflects the evaluator's training distribution, stylistic preferences, and inductive biases. The agent does not learn ``what works''---it learns ``what the evaluator likes.''

We term this phenomenon \textbf{Evaluator Preference Collapse (EPC)}: the progressive concentration of strategy weights toward strategies that satisfy the evaluator's inherent preferences, decoupled from objective task correctness.

\subsection{Contributions}

This paper makes four contributions:

\begin{enumerate}[leftmargin=*]
    \item \textbf{Identification of EPC.} We define and formalize Evaluator Preference Collapse as a systematic bias in self-referential agent optimization, and propose the Preference Collapse Index (PCI) as a quantitative metric.
    
    \item \textbf{Four-condition controlled experiment.} We compare same-model self-evaluation, cross-role evaluation, cross-model evaluation (GPT-4o-mini), and ground-truth benchmark verification. \textbf{All four conditions yield different optimal strategies}---zero agreement among evaluators.
    
    \item \textbf{Quantification of collapse magnitude.} Self-evaluation produces $84\%$ more strategy concentration than ground truth (PCI: $0.461$ vs $0.251$). The self-evaluated ``best'' strategy ranks only 5th of 8 under benchmark conditions.
    
    \item \textbf{Practical mitigations.} Cross-model evaluation reduces PCI by $17\%$; cross-role prompting reduces it by $2\%$ at zero additional API cost. We release our experimental framework for reproducible research.
\end{enumerate}

%===========================================================================
\section{Related Work}
%===========================================================================

\subsection{Test-Time Adaptation for LLMs}

Test-time training (TTT) was originally proposed for vision models adapting to distribution shifts at inference \cite{sun2020testtime}. Recent work has extended this concept to language agents. \citet{shinn2023reflexion} introduced Reflexion, where agents reflect on failures in natural language to improve subsequent attempts. \citet{madaan2023selfrefine} proposed Self-Refine, an iterative loop of generation and self-critique. \citet{yao2023tree} explored Tree-of-Thoughts, enabling exploration of multiple reasoning paths.

These methods share a common architecture: the agent generates output, evaluates it (using itself or a variant), and uses the evaluation to guide future behavior. None of them, however, systematically question whether the evaluation signal is unbiased. Our work is the first to demonstrate that self-evaluation bias can fundamentally redirect the agent's optimization trajectory.

\subsection{Self-Improving Language Models}

A parallel line of work focuses on LLMs that improve through self-generated data. SPIN \cite{chen2024spin} uses self-play fine-tuning where the model generates training data and learns from it iteratively. Self-Rewarding Language Models \cite{yuan2024selfrewarding} have the model act as its own reward model, creating an autonomous improvement loop. GRPO \cite{deepseek2024v3} compares outputs within groups for efficient reinforcement learning.

All these approaches implicitly assume that the model's self-evaluation is a reasonable proxy for true quality. Our findings challenge this assumption directly: the evaluator's preferences diverge systematically from objective task performance.

\subsection{LLM-as-Judge Evaluation}

A growing literature uses LLMs as evaluators for other LLMs \cite{zheng2023judging, li2024alpacaeval, chiang2024chatbot}. These works typically focus on inter-evaluator agreement and correlation with human judgments. However, they study \textit{static} evaluation scenarios---asking an LLM to rate outputs. Our work extends this to \textit{dynamic} settings where the evaluation signal drives an optimization process, revealing that small evaluator biases can compound over rounds of test-time learning into large strategy divergences.

\subsection{System-Level Evaluation in Production Models}

The GPT-4o System Card \cite{openai2024gpt4o} represents the most comprehensive public documentation of LLM evaluation methodology in a production system. It employs a multi-layered evaluation framework including automated benchmarks (MMLU, HumanEval), red-teaming, and human preference studies. Notably, the System Card acknowledges that ``safety alignment in multimodal scenarios is more complex'' and that ``multimodal hallucination risk increases'' with unified architectures. These observations are consistent with our EPC hypothesis: as evaluation scenarios become more complex and multimodal, the gap between evaluator preference and ground-truth correctness widens. Our work provides a quantitative framework (PCI) for measuring precisely this kind of evaluator-system divergence.

\subsection{Multi-Agent and Evaluator Diversity}

Recent multi-agent systems \cite{wu2023autogen, hong2023metagpt, chen2024chatdev} employ diverse agent roles, implicitly benefiting from evaluator diversity. Our work provides a theoretical grounding for this practice: evaluator diversity directly counteracts preference collapse. The Byzantine-resilient aggregation literature \cite{blanchard2017byzantine} offers complementary formal guarantees.

%===========================================================================
\section{Evaluator Preference Collapse}
%===========================================================================

\subsection{Problem Setup}

Consider an agent that maintains a set of $k$ reasoning strategies $\mathcal{S} = \{s_1, \ldots, s_k\}$, each associated with a weight $w_i \in (0,1)$ where $\sum_{i=1}^k w_i = 1$. At each round $t$, the agent:

\begin{enumerate}[leftmargin=*, label=(\arabic*)]
    \item Samples strategy $s \sim \text{Categorical}(\mathbf{w}^{(t)})$
    \item Executes $s$ on task $x$, producing output $o = \text{LLM}(s, x)$
    \item Evaluator $E$ assesses $o$, yielding reward $r = E(o, x) \in [0,1]$
    \item Updates weights: $w_s^{(t+1)} = w_s^{(t)} + \alpha \cdot (r - 0.5)$, renormalizes
\end{enumerate}

\subsection{The Collapse Mechanism}

The critical insight is in step (3). When the evaluator $E$ is the same model (or a model from the same family) as the executor, the reward $r$ is not $P(\text{correct} \mid o, x)$ but rather $P(\text{preferred by } E \mid o, x)$. Over $T$ rounds, the weight vector converges to:

\begin{equation}
\mathbf{w}^* = \arg\max_{\mathbf{w}} \mathbb{E}_{s \sim \mathbf{w}, x \sim \mathcal{D}} [E(\text{LLM}(s, x), x)]
\end{equation}

This differs from the \textit{task-optimal} weight vector:

\begin{equation}
\mathbf{w}^{\text{opt}} = \arg\max_{\mathbf{w}} \mathbb{E}_{s \sim \mathbf{w}, x \sim \mathcal{D}} [\text{TaskSuccess}(\text{LLM}(s, x), x)]
\end{equation}

The gap $\|\mathbf{w}^* - \mathbf{w}^{\text{opt}}\|$ quantifies the EPC effect---how far evaluator preference drives the agent from task-optimal behavior.

\subsection{Preference Collapse Index (PCI)}

We quantify EPC via the coefficient of variation of the weight vector:

\begin{equation}
\text{PCI}(\mathbf{w}) = \frac{\sigma(\mathbf{w})}{\mu(\mathbf{w})} = \frac{\sqrt{\frac{1}{k}\sum_{i=1}^k (w_i - \bar{w})^2}}{\bar{w}}
\end{equation}

\textbf{Properties:} PCI $= 0$ indicates perfectly uniform weights (no evaluator-driven concentration). PCI $> 0$ indicates progressive collapse toward preferred strategies. PCI is scale-normalized (divided by mean), making it comparable across different numbers of strategies $k$ and across experimental conditions. Higher PCI indicates stronger evaluator influence on strategy selection.

\subsection{Expected PCI Under Random Evaluation}

If the evaluator were unbiased (rewards independent of strategy), the expected weight after $T$ rounds would remain approximately uniform, yielding $\mathbb{E}[\text{PCI}] \approx 0$. Any significant deviation from zero---particularly systematic divergence across evaluator conditions---indicates evaluator-driven preference.

%===========================================================================
\section{Experimental Design}
%===========================================================================

\subsection{Four Evaluator Conditions}

We compare four evaluator configurations, varying the relationship between executor and evaluator:

\begin{table}[H]
\centering
\caption{Four experimental conditions varying evaluator configuration.}
\begin{tabular}{llll}
\toprule
\textbf{Condition} & \textbf{Executor} & \textbf{Evaluator} & \textbf{Key Property} \\
\midrule
A: Self-Eval & DeepSeek-chat & DeepSeek-chat (same role) & Standard TTRL baseline \\
B: Cross-Role & DeepSeek-chat & DeepSeek-chat (independent prompt) & Prompt-level mitigation \\
C: Benchmark & DeepSeek-chat & Ground-truth answer verification & Gold standard \\
D: Cross-Model & DeepSeek-chat & GPT-4o-mini (different family) & True cross-model evaluation \\
\bottomrule
\end{tabular}
\end{table}

\textbf{Condition A} represents the standard TTRL approach \cite{shinn2023reflexion, madaan2023selfrefine}. The evaluator prompt is minimal (``评估。更好？A或B：''), matching the executor's conversational style.

\textbf{Condition B} is our proposed lightweight mitigation. The evaluator receives a distinct system prompt: ``你是独立第三方AI评估员。客观公正,无任何偏好。只基于事实判断。'' This breaks role symmetry without changing the underlying model.

\textbf{Condition C} uses tasks with verifiably correct answers (mathematical, factual, programming). The reward is binary correctness rather than evaluator judgment. This serves as the ground-truth baseline.

\textbf{Condition D} uses GPT-4o-mini (accessed via api2d) as the evaluator for DeepSeek-chat outputs. This represents the strongest decoupling available without human evaluation.

\subsection{Tasks and Strategies}

\textbf{Tasks:} We use 15 tasks spanning five categories: (1) \textit{Verification}---6 tasks with verifiable correct answers (arithmetic, programming facts, Git commands); (2) \textit{Code Generation}---2 tasks requiring functional Python code; (3) \textit{Logical Reasoning}---2 tasks requiring multi-step deduction; (4) \textit{System Design}---2 tasks requiring architectural thinking; (5) \textit{Technical Explanation}---3 tasks requiring domain knowledge.

\textbf{Strategies:} We define $k=8$ reasoning strategies, each implemented as a system prompt prefix:

\begin{enumerate}[leftmargin=*,nosep]
    \item \texttt{step\_by\_step}: ``逐步推理,每步显式写出逻辑''
    \item \texttt{critical\_check}: ``先给出答案,再审视自身,指出漏洞并修正''
    \item \texttt{first\_principles}: ``从基本原理出发推导,不依赖记忆''
    \item \texttt{creative\_leap}: ``跳出常规思维,寻找创新方案''
    \item \texttt{analogy\_meta}: ``用类比和具体例子解释复杂概念''
    \item \texttt{evidence\_cite}: ``引用具体知识依据和论文证据''
    \item \texttt{synthesis}: ``综合多种视角,给出平衡全面的回答''
    \item \texttt{counterfactual}: ``考虑反面情况和边界条件''
\end{enumerate}

\subsection{Procedure}

For each condition, we run 10 rounds of TTRL. Each round consists of:

\begin{enumerate}[leftmargin=*, label=(\arabic*)]
    \item \textbf{Strategy selection}: Sample strategy $s$ from current weight distribution $\mathbf{w}$.
    \item \textbf{Execution}: Generate output $o = \text{DeepSeek}(s, \text{task})$.
    \item \textbf{Control generation}: Generate control output $c = \text{DeepSeek}(\texttt{step\_by\_step}, \text{task})$.
    \item \textbf{Evaluation}: Evaluator $E$ compares $o$ vs $c$, returning winner $\in \{\text{A}, \text{B}\}$.
    \item \textbf{Weight update}: $w_s \leftarrow w_s + 0.08$ if TTRL wins, $w_s \leftarrow w_s - 0.04$ otherwise. Weights are renormalized with minimum $0.01$.
\end{enumerate}

For Condition C (Benchmark), step (4) is replaced by direct verification against the known correct answer.

\subsection{Implementation Details}

All experiments use DeepSeek-chat accessed via the official API (\texttt{api.deepseek.com}) for execution. Cross-model evaluation uses GPT-4o-mini via the api2d proxy (\texttt{oa.api2d.net}). Temperature is set to $0.7$ for generation and $0.1$ for evaluation. Maximum tokens are $350$ for generation and $10$ for evaluation (single-letter output). Total experiment cost is approximately \$0.03 USD.

%===========================================================================
\section{Results}
%===========================================================================

\subsection{Preference Collapse Quantified}

\begin{table}[H]
\centering
\caption{PCI and top strategy for each evaluator condition. All four conditions produce different optimal strategies.}
\begin{tabular}{lccc}
\toprule
\textbf{Condition} & \textbf{PCI} & \textbf{Top Strategy} & \textbf{Top Weight} \\
\midrule
A: Self-Eval      & \textbf{0.461} & counterfactual  & 0.264 \\
B: Cross-Role     & \textbf{0.451} & synthesis       & 0.257 \\
C: Benchmark      & \textbf{0.251} & evidence\_cite  & 0.159 \\
D: GPT Cross-Model & \textbf{0.384} & creative\_leap  & 0.173 \\
\bottomrule
\end{tabular}
\end{table}

\textbf{Finding 1 (H1 confirmed):} Self-evaluation produces significantly more strategy concentration than ground-truth verification. The PCI ratio $\text{PCI}(A)/\text{PCI}(C) = 0.461/0.251 = 1.84$ represents an $84\%$ increase in preference collapse. Self-evaluation concentrates the top strategy at $26.4\%$ weight---more than double the uniform baseline of $12.5\%$.

\textbf{Finding 2 (H3 confirmed):} Cross-model evaluation partially mitigates collapse. $\text{PCI}(D) = 0.384 < \text{PCI}(A) = 0.461$, demonstrating a $17\%$ reduction. Using a different model family as evaluator decouples strategy evolution from executor bias.

\subsection{Zero Evaluator Agreement}

\begin{table}[H]
\centering
\caption{Complete strategy weight distributions across all four conditions. Bold indicates the optimal strategy for each condition. Underlining indicates strategies misranked by $\geq 3$ positions relative to ground truth.}
\small
\begin{tabular}{lcccc}
\toprule
\textbf{Strategy} & \textbf{A: Self} & \textbf{B: Cross-Role} & \textbf{C: Benchmark} & \textbf{D: GPT} \\
\midrule
counterfactual  & \textbf{0.264} & 0.047 & \underline{0.093} & 0.095 \\
synthesis       & 0.126 & \textbf{0.257} & 0.103 & 0.085 \\
evidence\_cite  & 0.064 & 0.134 & \textbf{0.159} & 0.068 \\
creative\_leap  & 0.149 & 0.058 & 0.111 & \textbf{0.173} \\
step\_by\_step  & 0.161 & 0.177 & 0.143 & 0.151 \\
critical\_check & 0.085 & 0.087 & 0.117 & 0.153 \\
first\_principles & 0.078 & 0.058 & 0.157 & 0.119 \\
analogy\_meta   & 0.073 & 0.183 & 0.113 & 0.156 \\
\bottomrule
\end{tabular}
\end{table}

\textbf{Finding 3 (H2 spectacularly confirmed):} All four evaluator conditions produce mutually different optimal strategies. There is \textit{zero agreement} among evaluators on what constitutes the best reasoning approach. The self-evaluated optimum (\texttt{counterfactual}, $26.4\%$) ranks only 5th of 8 under ground-truth conditions.

\textbf{Finding 4 (Misranking):} Self-evaluation systematically overestimates \texttt{counterfactual} (rank 1 vs. ground-truth rank 5) and \texttt{creative\_leap} (rank 3 vs. ground-truth rank 4), while underestimating \texttt{first\_principles} (rank 6 vs. ground-truth rank 2) and \texttt{evidence\_cite} (rank 7 vs. ground-truth rank 1). The Spearman rank correlation between self-evaluation and ground-truth rankings is $\rho = -0.024$---effectively random.

\subsection{PCI Gradient}

The monotonic decrease in PCI across evaluator conditions reveals a clear gradient:

\begin{equation}
\text{PCI: } \underbrace{0.461}_{\text{Self}} \;>\; \underbrace{0.451}_{\text{Cross-Role}} \;>\; \underbrace{0.384}_{\text{Cross-Model}} \;>\; \underbrace{0.251}_{\text{Ground Truth}}
\end{equation}

Each step away from the executor's own model family reduces evaluator preference collapse. The cross-role prompting alone provides a $2\%$ reduction (cost-free). Cross-model evaluation provides a $17\%$ reduction. Ground-truth anchoring provides the strongest debiasing at $46\%$ reduction from self-evaluation.

\subsection{Ablation: Strategy Win Rates}

Table~\ref{tab:wins} shows the per-strategy win rates against the fixed step\_by\_step control in pairwise A/B comparison. Strategies with positive win rates are reinforced; strategies with negative win rates are penalized over rounds.

\begin{table}[H]
\centering
\caption{Per-strategy win counts against fixed step\_by\_step control across 10 rounds per condition.}
\label{tab:wins}
\begin{tabular}{lcccc}
\toprule
\textbf{Strategy} & \textbf{A: Self} & \textbf{B: Cross} & \textbf{C: Bench} & \textbf{D: GPT} \\
\midrule
counterfactual  & 2 & 1 & N/A & 1 \\
synthesis       & 2 & 2 & N/A & 0 \\
creative\_leap  & 1 & 0 & N/A & 2 \\
step\_by\_step  & 0 & 1 & N/A & 1 \\
critical\_check & 0 & 1 & N/A & 2 \\
first\_principles & 1 & 0 & N/A & 1 \\
evidence\_cite  & 0 & 1 & N/A & 1 \\
analogy\_meta   & 0 & 0 & N/A & 0 \\
\bottomrule
\end{tabular}
\end{table}

Note that Condition C (Benchmark) uses direct answer verification rather than pairwise comparison, so win counts are not applicable. The sparse win distribution (2-3 wins maximum per strategy) indicates high variance in pairwise evaluation, consistent with known LLM-as-judge inconsistency \cite{zheng2023judging}.

%===========================================================================
\section{Discussion}
%===========================================================================

\subsection{Why EPC Matters for Agent Design}

Our findings challenge a fundamental assumption underlying the entire test-time agent self-improvement literature. If the evaluation signal is systematically biased by evaluator identity, then:

\begin{enumerate}[leftmargin=*]
    \item \textbf{Agents optimize for evaluator satisfaction, not task success.} The self-evaluated ``best'' strategy ranks 5th under ground truth. An agent that trusts its self-evaluation is optimizing for the wrong objective.
    
    \item \textbf{Reported self-improvement gains may be artifacts of evaluator bias.} If both the executor and evaluator share preferences, apparent improvement may reflect increasing alignment with shared biases rather than genuine capability gains.
    
    \item \textbf{Strategy diversity is valuable but fragile.} Self-evaluation concentrates weight on 1-2 strategies; ground truth distributes weight more evenly. This diversity loss makes agents brittle to task distribution shifts.
\end{enumerate}

\subsection{Theoretical Implications}

EPC can be understood through the lens of \textbf{reward hacking} in reinforcement learning \cite{amodei2016concrete}. When the reward signal is a proxy (evaluator judgment) rather than the true objective (task correctness), optimization finds strategies that maximize the proxy. Our contribution is demonstrating that even \textit{self-generated} reward signals---long assumed to be reasonably aligned with true quality---exhibit systematic proxy-reward divergence.

\subsection{Practical Recommendations}

Based on our findings, we recommend the following for agent self-improvement systems:

\begin{enumerate}[leftmargin=*, label=\textbf{R\arabic*}:]
    \item \textbf{Report PCI.} Include the Preference Collapse Index alongside accuracy metrics to quantify evaluator bias.
    \item \textbf{Use multi-evaluator baselines.} At minimum, include one cross-model evaluator and one ground-truth verifiable task set.
    \item \textbf{Diversify evaluator prompts.} Our cross-role condition demonstrates that even prompt-level changes reduce collapse.
    \item \textbf{Anchor with benchmark tasks.} Include verifiable tasks (math, code execution) to provide an objective signal that counterbalances evaluator bias.
\end{enumerate}

\subsection{Limitations and Future Work}

Our study has several limitations that suggest directions for future work. First, the 10-round experiment provides initial evidence but larger-scale replication (50+ rounds, 100+ tasks) would strengthen statistical conclusions. Second, our single task distribution may not represent all agent domains; EPC may manifest differently in code generation, dialogue, or embodied tasks. Third, GPT-4o-mini as a cross-model evaluator may itself exhibit evaluation biases inherited from its RLHF training---a phenomenon consistent with OpenAI's own observation in the GPT-4o System Card \cite{openai2024gpt4o} that ``safety alignment in multimodal scenarios is more complex'' and evaluator reliability degrades as task complexity increases. Fourth, we study only pairwise (A/B) comparison; EPC may differ under Likert-scale rating or ranking evaluation formats.

Future work should: (1) replicate EPC measurements across diverse model families (Claude, Gemini, Llama); (2) investigate whether larger models exhibit weaker or stronger EPC; (3) develop evaluator-agnostic reward signals using execution feedback, tool-use success, or human preference data; and (4) explore whether ensemble evaluators can cancel out individual evaluator biases.

%===========================================================================
\section{Conclusion}
%===========================================================================

We have identified and quantified \textbf{Evaluator Preference Collapse (EPC)}---a systematic bias in test-time agent self-improvement where strategy optimization converges to evaluator preferences rather than objective task performance. Through a four-condition controlled experiment, we demonstrated that: (1) self-evaluation produces $84\%$ more strategy concentration than ground truth; (2) all four evaluator conditions produce mutually different optimal strategies with zero agreement; and (3) cross-model evaluation and cross-role prompting offer practical, low-cost mitigations.

The broader implication is sobering for the agent self-improvement community: \textbf{your agent's ``optimal'' strategy depends on who is watching.} We urge the field to adopt multi-evaluator baselines, report PCI alongside standard metrics, and critically examine whether self-reported improvements reflect genuine capability gains or mere alignment with evaluator preferences. We release our experimental framework and all trajectory data to facilitate reproducible research in this direction.

%===========================================================================
\bibliographystyle{plain}
\begin{thebibliography}{99}

\bibitem{yao2023react}
S. Yao, J. Zhao, D. Yu, N. Du, I. Shafran, K. Narasimhan, and Y. Cao.
\textit{ReAct: Synergizing Reasoning and Acting in Language Models.}
ICLR, 2023.

\bibitem{wu2023autogen}
Q. Wu, G. Bansal, J. Zhang, Y. Wu, S. Zhang, E. Zhu, B. Li, L. Jiang, X. Zhang, and C. Wang.
\textit{AutoGen: Enabling Next-Gen LLM Applications via Multi-Agent Conversation.}
arXiv:2308.08155, 2023.

\bibitem{park2023generative}
J. S. Park, J. O'Brien, C. J. Cai, M. R. Morris, P. Liang, and M. S. Bernstein.
\textit{Generative Agents: Interactive Simulacra of Human Behavior.}
UIST, 2023.

\bibitem{shinn2023reflexion}
N. Shinn, F. Cassano, A. Gopinath, K. Narasimhan, and S. Yao.
\textit{Reflexion: Language Agents with Verbal Reinforcement Learning.}
NeurIPS, 2023.

\bibitem{madaan2023selfrefine}
A. Madaan, N. Tandon, P. Gupta, S. Hallinan, L. Gao, S. Wiegreffe, U. Alon, N. Dziri, S. Prabhumoye, Y. Yang, et al.
\textit{Self-Refine: Iterative Refinement with Self-Feedback.}
NeurIPS, 2023.

\bibitem{sun2020testtime}
Y. Sun, X. Wang, Z. Liu, J. Miller, A. A. Efros, and M. Hardt.
\textit{Test-Time Training with Self-Supervision for Generalization under Distribution Shifts.}
ICLR, 2020.

\bibitem{yuan2024selfrewarding}
W. Yuan, R. Y. Pang, K. Cho, S. Sukhbaatar, J. Xu, and J. Weston.
\textit{Self-Rewarding Language Models.}
ICML, 2024.

\bibitem{chen2024spin}
Z. Chen, Y. Deng, H. Yuan, K. Ji, and Q. Gu.
\textit{SPIN: Self-Play Fine-Tuning Converts Weak LLMs to Strong LLMs.}
arXiv:2401.01335, 2024.

\bibitem{deepseek2024v3}
DeepSeek-AI.
\textit{DeepSeek-V3 Technical Report.}
arXiv:2412.19437, 2024.

\bibitem{yao2023tree}
S. Yao, D. Yu, J. Zhao, I. Shafran, T. L. Griffiths, Y. Cao, and K. Narasimhan.
\textit{Tree of Thoughts: Deliberate Problem Solving with Large Language Models.}
NeurIPS, 2023.

\bibitem{hong2023metagpt}
S. Hong, X. Zheng, J. Chen, Y. Cheng, J. Wang, C. Zhang, Z. Wang, S. K. S. Yau, Z. Lin, L. Zhou, et al.
\textit{MetaGPT: Meta Programming for Multi-Agent Collaborative Framework.}
arXiv:2308.00352, 2023.

\bibitem{chen2024chatdev}
Q. Chen, X. Du, C. Liu, J. Wang, Z. Wang, and X. Zhang.
\textit{ChatDev: Communicative Agents for Software Development.}
ACL, 2024.

\bibitem{zheng2023judging}
L. Zheng, W.-L. Chiang, Y. Sheng, S. Zhuang, Z. Wu, Y. Zhuang, Z. Lin, Z. Li, D. Li, E. P. Xing, et al.
\textit{Judging LLM-as-a-Judge with MT-Bench and Chatbot Arena.}
NeurIPS, 2023.

\bibitem{li2024alpacaeval}
X. Li, T. Zhang, Y. Dubois, R. Taori, I. Gulrajani, C. Guestrin, P. Liang, and T. B. Hashimoto.
\textit{AlpacaEval: An Automatic Evaluator of Instruction-following Models.}
ICLR, 2024.

\bibitem{chiang2024chatbot}
W.-L. Chiang, L. Zheng, Y. Sheng, A. N. Angelopoulos, T. Li, D. Li, H. Zhang, B. Zhu, M. Jordan, J. E. Gonzalez, and I. Stoica.
\textit{Chatbot Arena: An Open Platform for Evaluating LLMs by Human Preference.}
ICML, 2024.

\bibitem{blanchard2017byzantine}
P. Blanchard, E. M. El Mhamdi, R. Guerraoui, and J. Stainer.
\textit{Machine Learning with Adversaries: Byzantine Tolerant Gradient Descent.}
NeurIPS, 2017.

\bibitem{openai2024gpt4o}
OpenAI.
\textit{GPT-4o System Card.}
arXiv:2410.21276, 2024.

\bibitem{amodei2016concrete}
D. Amodei, C. Olah, J. Steinhardt, P. Christiano, J. Schulman, and D. Man\'e.
\textit{Concrete Problems in AI Safety.}
arXiv:1606.06565, 2016.

\end{thebibliography}

%===========================================================================
\appendix
\section{Experiment Reproducibility}
%===========================================================================

\subsection{Requirements}

Python 3.8+, DeepSeek API key, no GPU required. All experiments run via API calls.

\subsection{Running the Experiment}

\begin{verbatim}
export DEEPSEEK_API_KEY="sk-..."
python epc_experiment.py
\end{verbatim}

\subsection{Output}

The script produces \texttt{experiments/epc\_v2\_4conditions.json} containing per-round strategy weights, PCI values, win/loss records, and correctness rates for all four conditions. The file structure is:

\begin{verbatim}
{
  "A_SELF": {"top": "...", "pci": 0.461, "weights": {...}},
  "B_CROSS": {"top": "...", "pci": 0.451, "weights": {...}},
  "C_BENCH": {"top": "...", "pci": 0.251, "weights": {...}},
  "D_GPT": {"top": "...", "pci": 0.384, "weights": {...}}
}
\end{verbatim}

\subsection{Repository Structure}

\begin{verbatim}
aettl-research/
  epc_experiment.py           # Main experiment
  paper/epc_paper.tex          # This paper (LaTeX)
  experiments/
    epc_v2_4conditions.json    # Raw data
    papers_90.json             # Paper corpus
    reflection_90papers.json   # Cross-domain synthesis
\end{verbatim}

\end{document}
